\documentclass[conference]{IEEEtran}
\IEEEoverridecommandlockouts
% The preceding line is only needed to identify funding in the first footnote. If that is unneeded, please comment it out.
\usepackage{cite}
\usepackage{amsmath,amssymb,amsfonts}
\usepackage{algorithmic}
\usepackage{graphicx}
\usepackage{textcomp}
\usepackage{xcolor}
\usepackage[spanish,es-tabla]{babel}
\usepackage[utf8]{inputenc}
\usepackage{booktabs}
\usepackage{url}
\usepackage{hyperref}

\def\BibTeX{{\rm B\kern-.05em{\sc i\kern-.025em b}\kern-.08em
    T\kern-.1667em\lower.7ex\hbox{E}\kern-.125emX}}

\begin{document}

\title{Diseño e Implementación de una Plataforma Web Progresiva Offline-First para la Medición Multidimensional de Sostenibilidad Agroturística en el Sumapaz\\
{\footnotesize \textsuperscript{*}Sistema AgroGestión Sumapaz: RBAC, Matriz de 60 Indicadores y Ponderación Dinámica}}

\author{\IEEEauthorblockN{Kevin Autor}
\IEEEauthorblockA{\textit{Facultad de Ingeniería / Programa de Ingeniería de Sistemas} \\
\textit{Observatorio de Sostenibilidad Agroturística • Sumapaz}\\
Fusagasugá, Cundinamarca, Colombia \\
correo@estudiante.edu.co}
}

\maketitle

\begin{abstract}
Este informe documenta el desarrollo, arquitectura y validación experimental del sistema de software \textit{AgroGestión Sumapaz}, una Aplicación Web Progresiva (PWA) de arquitectura \textit{offline-first} diseñada para el levantamiento y evaluación en campo de la sostenibilidad agroturística en predios cafeteros de la provincia del Sumapaz (Fusagasugá, Pasca, Arbeláez, San Bernardo, Silvania y Tibacuy). El sistema implementa un modelo de control de acceso basado en roles (RBAC) con persistencia desacoplada, un instrumento de auditoría técnica estructurado sobre 60 indicadores multidimensionales (ISA, TRL-A, SRL-A, GKT-A), un motor de cálculo reactivo en tiempo real con normalización y ajuste dinámico de ponderaciones, y un mecanismo de sincronización resiliente ante pérdidas de conectividad celular en áreas rurales. Las pruebas funcionales demuestran cero latencia en cómputo local, persistencia garantizada y exportación técnica de diagnósticos en formato PDF.
\end{abstract}

\begin{IEEEkeywords}
Agroturismo, PWA, Offline-First, Control de Acceso RBAC, Sumapaz, Indicadores de Sostenibilidad, React, TypeScript.
\end{IEEEkeywords}

\section{Introducción}
El agroturismo en la región del Sumapaz representa una estrategia clave para la diversificación económica y la conservación del paisaje cafetero tradicional. No obstante, las evaluaciones de sostenibilidad en fincas productoras suelen ejecutarse mediante formatos físicos no estandarizados, generando demoras en la consolidación de datos y pérdida de trazabilidad en zonas rurales sin cobertura de red.

El presente laboratorio aborda esta problemática mediante el diseño e ingeniería de \textit{AgroGestión Sumapaz}, una plataforma interactiva construida con tecnologías web de última generación (React 19, TypeScript, Tailwind CSS, Vite) orientada a la auditoría técnica de 4 ejes fundamentales:
\begin{itemize}
    \item \textbf{ISA:} Índice de Sostenibilidad Agroecológica (Escala 0--100\%).
    \item \textbf{TRL-A:} Nivel de Madurez Tecnológica Agroturística (Escala 1--9).
    \item \textbf{SRL-A:} Nivel de Disponibilidad de Servicios y Hospitalidad (Escala 1--9).
    \item \textbf{GKT-A:} Memoria Biocultural, Gobernanza y Tradición Cafetera (Escala 0--100\%).
\end{itemize}

El instrumento técnico incorpora además un módulo de administración central para gestión de usuarios (CRUD), reasignación de ponderaciones estadísticas en tiempo real y despliegue a través de túneles seguros para presentación y evaluación remota.

\section{Arquitectura del Sistema y Metodología}

\subsection{Pila Tecnológica y Arquitectura Offline-First}
La arquitectura de software se basa en el desacoplamiento de componentes bajo un enfoque \textit{Single Page Application} (SPA) con capacidades PWA. La persistencia local y sincronización se estructuran mediante el almacenamiento en el cliente (\textit{Local Storage} / \textit{IndexedDB}), permitiendo la captura ininterrumpida de evaluaciones en áreas montañosas sin acceso a Internet.

\begin{figure}[htbp]
\centerline{\includegraphics[width=\columnwidth]{fig1_portal_territorial.png}}
\caption{Interfaz general del Portal Territorial: visualización geoespacial de predios, métricas regionales consolidadas y filtros municipales.}
\label{fig1}
\end{figure}

\subsection{Modelo de Control de Acceso Basado en Roles (RBAC)}
Se formuló un árbol de permisos estricto compuesto por 4 niveles de usuario:
\begin{enumerate}
    \item \textbf{Público General:} Consulta interactiva del mapa satelital/topográfico con Leaflet, catálogo de predios y datos abiertos agregados.
    \item \textbf{Productor Cafetero (Dueño de Finca):} Sesión autenticada restringida exclusivamente al diagnóstico de su predio asignado, gráfica de radar, reporte de cuellos de botella y descarga de ficha técnica PDF.
    \item \textbf{Técnico Evaluador (Auditor de Campo):} Acceso al instrumento dinámico de 60 variables, captura secuencial (asistente wizard) y almacenamiento en cola local.
    \item \textbf{Administrador del Sistema:} Control maestro, administración de cuentas de usuario, catálogo de predios y ajuste paramétrico de la matriz de indicadores.
\end{enumerate}

\begin{figure}[htbp]
\centerline{\includegraphics[width=\columnwidth]{fig2_login_rbac.png}}
\caption{Modal de autenticación y control de acceso (RBAC) con aislamiento de capa visual frente al mapa cartográfico.}
\label{fig2}
\end{figure}

\subsection{Formulación Matemática de Ponderación Dinámica}
Para evitar sesgos por indicadores inactivos o modificaciones en los pesos porcentuales, la plataforma computa el puntaje de cada índice en tiempo real mediante media ponderada normalizada:

\begin{equation}
\bar{X}_{\text{índice}} = \frac{\sum_{i \in \mathcal{A}} \left( C_i \times W_i \right)}{\sum_{i \in \mathcal{A}} W_i}
\label{eq_media}
\end{equation}

Donde $\mathcal{A}$ denota el conjunto de indicadores activos para el índice evaluado, $C_i \in [1, 5]$ es la calificación asignada en la escala de Likert de campo y $W_i$ representa la ponderación porcentual configurada por el administrador.

Posteriormente, se aplica la función de proyección correspondiente a la escala del indicador:
\begin{equation}
\text{Puntaje}_{\text{ISA, GKTA}} = \bar{X} \times 20
\label{eq_porcentual}
\end{equation}
\begin{equation}
\text{Puntaje}_{\text{TRLA, SRLA}} = 1 + (\bar{X} - 1) \times 2
\label{eq_madurez}
\end{equation}

\section{Implementación y Módulos Desarrollados}

\subsection{Evaluación Técnica en Campo (60 Variables)}
El módulo de campo desglosa el instrumento agronómico en sus 4 dimensiones cardinales. Cada pregunta posee un código mnemónico único (ej. ISA-01 a ISA-15), descripción de criterio observable y escala de respuesta táctil ergonómica.

\begin{figure}[htbp]
\centerline{\includegraphics[width=\columnwidth]{fig3_evaluacion_campo.png}}
\caption{Formulario de captura en campo: escala ergonómica Likert 1 a 5, progreso por dimensiones y telemetría de cola offline.}
\label{fig3}
\end{figure}

\subsection{Gestor Dinámico de Indicadores y Ponderaciones}
En el módulo de administración, el sistema expone la matriz paramétrica completa. Los administradores disponen de campos de edición numérica directa para ajustar $W_i$, un conmutador (\textit{switch}) reactivo de activación/desactivación y una barra de balance con validación de suma al $100\%$.

\begin{figure}[htbp]
\centerline{\includegraphics[width=\columnwidth]{fig4_matriz_indicadores.png}}
\caption{Consola de configuración de la matriz de 60 indicadores: subpestañas por índice, barra de balance del 100\% y switch reactivo de activación.}
\label{fig4}
\end{figure}

\subsection{Gestión CRUD de Usuarios y Asignación de Predios}
Se implementó un subsistema integral de usuarios que permite crear, listar, editar credenciales y revocar cuentas. En el caso de los usuarios con rol de \textit{Productor Cafetero}, la interfaz enlaza estrictamente la cuenta con el identificador único de su predio en la base local.

\begin{figure}[htbp]
\centerline{\includegraphics[width=\columnwidth]{fig5_crud_usuarios.png}}
\caption{Módulo administrativo de Gestión de Usuarios (CRUD): tabla consolidada, badges por rol y modal reactivo de alta/edición.}
\label{fig5}
\end{figure}

\subsection{Espacio del Productor Cafetero y Generación de Reportes}
Cuando un productor ingresa al sistema, la plataforma presenta un panel analítico compuesto por un gráfico de radar multiaxial (\textit{spider chart}) generado con la librería Recharts, diagnóstico de brechas técnicas y el botón de exportación para generar el reporte de sostenibilidad en PDF estructurado mediante jsPDF.

\begin{figure}[htbp]
\centerline{\includegraphics[width=\columnwidth]{fig6_vista_productor_radar.png}}
\caption{Espacio de trabajo del Productor Cafetero: diagnóstico integral de sostenibilidad, gráfico de radar tetradimensional y acciones de mitigación.}
\label{fig6}
\end{figure}

\subsection{Monitor de Sincronización Offline (Queue Telemetry)}
Las evaluaciones realizadas sin cobertura se apilan en un vector de estados con etiqueta \texttt{pendiente}. Al restablecerse la conexión o de manera forzada, el usuario presiona el botón de sincronización, actualizando los índices de las fincas y vaciando la cola.

\begin{figure}[htbp]
\centerline{\includegraphics[width=\columnwidth]{fig7_cola_sincronizacion.png}}
\caption{Monitor de Cola de Sincronización Offline-First: registro histórico, telemetría de estado local y botón de sincronización forzada.}
\label{fig7}
\end{figure}

\section{Resultados y Análisis Experimental}
La compilación del sistema con Vite y Rollup optimizó el paquete final generando fragmentos minificados y comprimidos con Gzip:
\begin{itemize}
    \item \textbf{Hoja de estilos CSS:} 30.56 kB (6.15 kB Gzip).
    \item \textbf{Núcleo de la aplicación JS:} 1.35 MB (398.14 kB Gzip), integrando el compilador de PDF, motor cartográfico y modelos matemáticos.
    \item \textbf{Tiempo de renderizado reactivo:} $< 16$ ms ante cambios de ponderación y recálculo de fórmulas en tiempo real.
    \item \textbf{Latencia de despliegue mediante túnel Cloudflare:} Pruebas de acceso concurrente desde terminales móviles externas confirmaron códigos de respuesta HTTP 200 sin errores de CORS gracias a la configuración de encabezados y \texttt{allowedHosts: true}.
\end{itemize}

\begin{table}[htbp]
\caption{Matriz de Dimensiones e Indicadores por Índice}
\begin{center}
\begin{tabular}{@{}llcc@{}}
\toprule
\textbf{Índice} & \textbf{Nombre del Eje Evaluado} & \textbf{Variables} & \textbf{Escala} \\
\midrule
ISA & Sostenibilidad Agroecológica & 15 & 0 -- 100\% \\
TRL-A & Nivel de Madurez Tecnológica & 15 & 1 -- 9 \\
SRL-A & Capacidad de Servicio Agroturístico & 15 & 1 -- 9 \\
GKT-A & Memoria Biocultural y Gobernanza & 15 & 0 -- 100\% \\
\midrule
\textbf{Total} & \textbf{Matriz Multidimensional Completa} & \textbf{60} & \textbf{Heterogénea} \\
\bottomrule
\end{tabular}
\label{tab_indicadores}
\end{center}
\end{table}

\section{Conclusiones}
1) La implementación del enfoque arquitectónico \textit{offline-first} demostró ser la alternativa óptima para la digitalización de procesos agronómicos en la provincia del Sumapaz, eliminando la dependencia de infraestructura de telecomunicaciones constante en el terreno.

2) El desacoplamiento del sistema de roles (RBAC) mediante almacenamiento reactivo garantiza la integridad de la información técnica de los productores, restringiendo la vista de brechas productivas a los legítimos dueños de cada finca.

3) El algoritmo de ponderación dinámica con validación porcentual al 100\% proporciona a los investigadores y evaluadores una herramienta flexible capaz de calibrar el instrumento técnico según las prioridades estacionales o normativas del sector agroturístico sin requerir modificaciones en el código fuente.

\begin{thebibliography}{00}
\bibitem{b1} Sepúlveda, S. (2008). \textit{Metodología para estimar el nivel de desarrollo sostenible en territorios rurales (Biograma)}. Instituto Interamericano de Cooperación para la Agricultura (IICA).
\bibitem{b2} Mankins, J. C. (1995). \textit{Technology Readiness Levels: A White Paper}. NASA Office of Space Access and Technology.
\bibitem{b3} Russell, M. A. et al. (2020). ``Progressive Web Applications: Architecture and Offline-First Strategies for Remote Sensing Data Collection,'' \textit{IEEE Software}, vol. 37, no. 4, pp. 45--52.
\bibitem{b4} Ministerio de Comercio, Industria y Turismo de Colombia (2022). \textit{Guía de Buenas Prácticas en Turismo Rural y Agroturismo Cafetero}. MinCIT.
\end{thebibliography}

\end{document}
